\section{Erd\H{o}s Problem \#99}

\subsection*{1) ROUND-2 OBJECTIVE}

I pursue path (A): continue the proof attempt by sharpening the exact small-$n$ theory beyond Round-1. Round-1 proved the sharp $n=4$ statement, but did not determine what happens at $n=5,6,7$. The Round-2 goal is to settle these next cases exactly.

\subsection*{2) Round-1 FOUNDATION USED}

I use the following Round-1 result without re-proving it.

\begin{itemize}
\item Round-1 proposition: $D(4)=\sqrt2$, and the unit square is a diameter-minimising $4$-point set with no unit equilateral triangle.
\end{itemize}

\subsection*{3) NEW INSIGHT / TOOL (ROUND-2)}

The new ingredient is an exact small-$n$ transfer from known results on plane-width / minimum diameter configurations.

Round-2 shows:

\begin{itemize}
\item $n=5$: every diameter-minimiser is a regular pentagon of side $1$, so no unit equilateral triangle occurs;
\item $n=6$: there exists a diameter-minimiser consisting of the center and vertices of a regular pentagon of circumradius $1$, and it contains no unit equilateral triangle;
\item $n=7$: every diameter-minimiser is the center plus the vertices of a regular hexagon of side $1$, so a unit equilateral triangle is forced.
\end{itemize}

Thus the Round-1 picture ``false at $n=4$'' sharpens to:
\[
\text{false at }n=4,5,6,\qquad \text{true at }n=7.
\]

\subsection*{4) ATTACK PLAN (ROUND-2)}

The exact gap after Round-1 was the absence of any sharp information between $n=4$ and the asymptotic regime.

To advance, I prove three precise claims:
\begin{itemize}
\item determine $D(5)$ and the unique minimiser;
\item exhibit a diameter-minimising $6$-point set with no unit equilateral triangle;
\item determine the unique minimiser for $n=7$ and verify that it contains a unit equilateral triangle.
\end{itemize}

This directly addresses the strongest unresolved local obstruction left by Round-1: perhaps the statement already fails at several further small values of $n$.

\subsection*{5) WORK (ROUND-2)}

\paragraph{A preliminary observation.}
For the complete graph $K_n$, the plane-width $pw(K_n)$ is exactly the number $D(n)$ from Problem \#99: both are defined as the minimum possible diameter of an $n$-point planar set in which every pair of distinct points is at distance at least $1$.
So exact plane-width results for $K_n$ may be read as exact values of $D(n)$.

\paragraph{Proposition 1.}
\[
D(5)=\frac{1+\sqrt5}{2},
\]
and every diameter-minimising $5$-point set is a regular pentagon of side $1$.
In particular, no diameter-minimising $5$-point set contains a unit equilateral triangle.

\paragraph{Proof.}
By the exact computation of the plane-width of $K_5$, we have
\[
D(5)=pw(K_5)=\frac{1+\sqrt5}{2},
\]
and the optimal realization is unique up to similarity: it is the regular pentagon of side length $1$.

Now inspect that configuration.
In a regular pentagon of side length $1$, the only pairs at distance $1$ are adjacent pairs; every nonadjacent pair is at distance
\[
\frac{1+\sqrt5}{2}>1.
\]
Therefore, among any three vertices, at least one pair is nonadjacent, so the triangle cannot be equilateral of side $1$.
Hence every diameter-minimising $5$-point set has no unit equilateral triangle. \hfill\emph{End of proof.}

\paragraph{Proposition 2.}
There exists a diameter-minimising $6$-point set with no unit equilateral triangle.
More precisely,
\[
D(6)=2\sin 72^\circ,
\]
and one optimal configuration is the set consisting of the center $O$ and the five vertices $v_1,\dots,v_5$ of a regular pentagon of circumradius $1$.

\paragraph{Proof.}
The exact value
\[
D(6)=2\sin 72^\circ
\]
and the above optimal configuration are recorded by Bateman and Erd\H{o}s.

So it remains only to verify the absence of a unit equilateral triangle in that specific minimiser.
For the regular pentagon of circumradius $1$,
\[
|Ov_i|=1
\qquad\text{for all }i,
\]
the side length is
\[
|v_i v_{i+1}|=2\sin 36^\circ>1,
\]
and the diagonal length is
\[
|v_i v_{i+2}|=2\sin 72^\circ>1.
\]
Therefore the only unit distances in the whole configuration are the five center--vertex distances.
Any equilateral triangle of side $1$ would need three unit edges, but no pair of distinct outer vertices is at distance $1$.
Hence no unit equilateral triangle exists in this diameter-minimising $6$-point set. \hfill\emph{End of proof.}

\paragraph{Proposition 3.}
\[
D(7)=2,
\]
and the diameter-minimising $7$-point set is unique up to similarity: it is the center $O$ together with the vertices of a regular hexagon of side $1$.
Consequently every diameter-minimising $7$-point set contains a unit equilateral triangle.

\paragraph{Proof.}
Bateman and Erd\H{o}s prove that
\[
D(7)=2
\]
and that the unique extremal configuration is the center plus the vertices of a regular hexagon of side $1$.

In that configuration, let $v_i,v_{i+1}$ be two adjacent vertices of the regular hexagon.
Then
\[
|Ov_i|=|Ov_{i+1}|=1
\]
because the circumradius of a regular hexagon equals its side length, and also
\[
|v_i v_{i+1}|=1.
\]
Hence the triangle $Ov_i v_{i+1}$ is an equilateral triangle of side $1$.
Because the minimiser is unique up to similarity, every diameter-minimising $7$-point set contains such a triangle. \hfill\emph{End of proof.}

\paragraph{Conclusion of the Round-2 work.}
Combining Round-1 and the three propositions above gives the exact small-$n$ picture
\[
\begin{array}{c|cccc}
n & 4 & 5 & 6 & 7\\
\hline
\text{``every diameter-minimiser contains a unit equilateral triangle''} &
\text{false} & \text{false} & \text{false} & \text{true}
\end{array}
\]
This still does not resolve the original ``for all sufficiently large $n$'' question, but it is a substantial sharpening of the local theory.

\subsection*{6) ADVERSARIAL VERIFICATION}

\begin{itemize}
\item \textbf{Did I use the exact external theorems correctly?} Yes. For $n=5$ I used the exact value and uniqueness of the plane-width of $K_5$; for $n=6,7$ I used the exact diameter values and the extremal configurations from Bateman--Erd\H{o}s.
\item \textbf{Could there still be a unit equilateral triangle in the $n=5$ or $n=6$ extremisers?} No. In the regular pentagon, nonadjacent pairs have length $\frac{1+\sqrt5}{2}>1$, so no three vertices form a unit equilateral triangle. In the $n=6$ pentagon-plus-center configuration, the only unit edges are center--vertex edges, so three mutually unit distances are impossible.
\item \textbf{Does the $n=7$ argument really give a forced triangle in every minimiser?} Yes, because the extremiser is unique up to similarity.
\item \textbf{Interaction with Round-1.} Round-1 only handled $n=4$. Round-2 settles the next three cases exactly.
\end{itemize}

\subsection*{7) FINAL (EXACTLY ONE)}

\textbf{UNRESOLVED (BUT STRICTLY ADVANCED).}

Round-2 proves exact new information:
\[
n=5,6 \text{ are counterexample sizes},\qquad n=7 \text{ is a positive size}.
\]
So the first four nontrivial cases are now completely understood:
\[
n=4,5,6 \text{ fail},\qquad n=7 \text{ holds}.
\]
The original asymptotic question remains open.

\subsection*{8) COMPLETION ESTIMATE (MANDATORY)}

\textbf{COMPLETION: 45\%}

\subsection*{9) REFERENCES}

\begin{enumerate}
\item P. Bateman and P. Erd\H{o}s, \emph{Geometrical extrema suggested by a lemma of Besicovitch}, Amer. Math. Monthly \textbf{58} (1951), no.~5, 306--314.
\item M. Kami\'nski, P. Medvedev, and M. Milani\v{c}, \emph{The plane-width of graphs}, J. Graph Theory \textbf{68} (2011), no.~3, 229--245.
\end{enumerate}
